% =====================================================================
%  Projektdokumentation
%  Phase 1: Projektübersicht, Risikomanagement, Zeitplanung
%  Phase 2: + Link zum GitHub-Repository (Versionsverwaltung)
%  Phase 3: + Benutzeranleitung (Installation, Ausführung),
%           + Anforderungserfüllung / technische Schulden
%  Quellcode-Anhang (maschinenlesbar) ab Phase 2 automatisch.
% =====================================================================
\documentclass{iu-pse}
\input{meta}
\addbibresource{literatur.bib}

\dokumenttitel{Projektdokumentation}
\dokumentuntertitel{\Projektname}

\begin{document}
\titelblatt
\verzeichnisse
% Pflicht ab drei Abbildungen bzw. drei Tabellen (Richtlinien Kap. 3.1):
\abbildungsverzeichnis
\tabellenverzeichnis

\begin{abkuerzungen}[MoSCoW]
  \acro{GUI}{Graphical User Interface}
  \acro{MVP}{Minimum Viable Product}
  \acro{MoSCoW}{Must, Should, Could, Won't}
\end{abkuerzungen}

% \genderhinweis   % optional

\hauptteil
% =====================================================================
%  1 Projektübersicht  (Aufgabenstellung 1.1: „Beschreibung des Projekts
%  und des angestrebten Ergebnisses“)
%  Leitfragen Prüfungsleitfaden Phase 1:
%   - Was ist das Ziel meiner Arbeit?  - Was ist meine Idee?
%   - Ideenprüfung: Was spricht für, was gegen die Idee?
%   - Was ist das Konzept?  - Wie plane ich vorzugehen?
%   - Welche Methodik/welches Tool habe ich gewählt und warum?
%  Bewertung „Problemabgrenzung/Zielsetzung“ (10 %): Ziele, Umfang,
%  angestrebtes Ergebnis klar; Nutzen für die Zielgruppe deutlich.
%  Hinweis: Kein „ich“, „wir“, „man“ (Richtlinien Kap. 1).
%  Jedes (Unter-)Kapitel mind. ½ Seite; Unterteilung nur mit ≥ 2 Unterpunkten.
% =====================================================================
\section{Projektübersicht}
\TODO{2--3 Sätze Hinführung zu den Unterkapiteln}

\subsection{Ausgangssituation und Motivation}
\TODO{Problem, das die Work-Life-Balance-App löst; Kontext; Relevanz mit Quellen belegen}

\subsection{Zielsetzung und Kundennutzen}
\TODO{Ziel des Projekts; Zielgruppe (Details im Anforderungsdokument); klar definierter Customer Value}

\subsection{Projektumfang und Abgrenzung}
\TODO{Was ist Teil des Projekts (Kernfunktionen), was ausdrücklich nicht? Abgrenzung zu bestehenden Apps (Kreativität, Neuheit); Ideenprüfung: Was spricht dafür, was dagegen?}

\subsection{Angestrebtes Ergebnis}
\TODO{Funktionsfähige Windows-Desktopanwendung in Python mit GUI und Testdaten, Dokumentation, GitHub-Repository}

\subsection{Vorgehensmodell und Werkzeuge}
\TODO{Gewähltes Vorgehen (z. B. iterativ-inkrementell) mit Begründung und Quelle; Werkzeuge: LaTeX, draw.io, Git/GitHub, Python}

% =====================================================================
%  2 Risikomanagement  (Aufgabenstellung 1.1: „Identifikation potenzieller
%  Projektrisiken und Gegenmaßnahmen inkl. Bewertung der Risiken nach
%  Eintrittswahrscheinlichkeit und Schwere der potenziellen Auswirkungen“)
% =====================================================================
\section{Risikomanagement}
\TODO{Kurze Hinführung}

\subsection{Methodik der Risikobewertung}
\TODO{Skalen erläutern (z. B. Eintrittswahrscheinlichkeit W und Schwere S jeweils 1--5; Risikowert R = W $\times$ S); Schwellenwerte für Handlungsbedarf; Quelle}

\subsection{Risikoliste und Gegenmaßnahmen}
\TODO{Risiken ergänzen (technisch, organisatorisch, zeitlich, fachlich); Beispielzeile ersetzen}

\begin{longtable}{@{}L{1.0cm}L{4.4cm}>{\centering\arraybackslash}p{0.6cm}>{\centering\arraybackslash}p{0.6cm}>{\centering\arraybackslash}p{0.7cm}L{6.4cm}@{}}
  \caption{Risikoliste mit Bewertung und Gegenmaßnahmen}\label{tab:risiken}\\
  \toprule
  \textbf{ID} & \textbf{Risiko} & \textbf{W} & \textbf{S} & \textbf{R} & \textbf{Gegenmaßnahme} \\
  \midrule
  \endfirsthead
  \toprule
  \textbf{ID} & \textbf{Risiko} & \textbf{W} & \textbf{S} & \textbf{R} & \textbf{Gegenmaßnahme} \\
  \midrule
  \endhead
  \bottomrule
  \endfoot
  R01 & \TODO{Beispiel: Zeitplan wird durch berufliche Belastung überschritten} & 3 & 4 & 12 & \TODO{Puffer je Meilenstein, Priorisierung nach MoSCoW} \\
  R02 & \TODO{} & & & & \TODO{} \\
\end{longtable}
\quelle{Eigene Darstellung.}

\subsection{Risikomatrix}
\TODO{Risiken in einer Matrix (W über S) einordnen; Grafik als docs/diagramme/export/risikomatrix.pdf}
\begin{figure}[H]
  \caption{Risikomatrix}\label{abb:risikomatrix}
  \centering
  \grafik[0.7\linewidth]{risikomatrix}
  \quelle{Eigene Darstellung.}
\end{figure}

% =====================================================================
%  3 Zeitplanung  (Aufgabenstellung 1.1: „hinreichend detaillierter
%  Zeitplan und Erfassung wesentlicher Arbeitspakete und Meilensteine“)
%  ACHTUNG Bewertung: „Wurde der initial festgelegte Zeitplan eingehalten?“
%  -> realistisch planen; Abweichungen später im Abstract reflektieren.
%  Feedback der Tutor:innen kann bis zu 7 Tage dauern -> einplanen!
% =====================================================================
\section{Zeitplanung}
\TODO{Kurze Hinführung; Planungszeitraum; verfügbare Wochenstunden}

\subsection{Arbeitspakete}
\begin{longtable}{@{}L{1.2cm}L{6.3cm}L{1.6cm}L{1.4cm}L{4.6cm}@{}}
  \caption{Arbeitspakete}\label{tab:arbeitspakete}\\
  \toprule
  \textbf{AP} & \textbf{Arbeitspaket} & \textbf{Phase} & \textbf{Std.} & \textbf{Ergebnis} \\
  \midrule
  \endfirsthead
  \toprule
  \textbf{AP} & \textbf{Arbeitspaket} & \textbf{Phase} & \textbf{Std.} & \textbf{Ergebnis} \\
  \midrule
  \endhead
  \bottomrule
  \endfoot
  AP1  & Projektplanung & 1 & \TODO{} & Projektdokumentation \\
  AP2  & Anforderungsermittlung & 1 & \TODO{} & Anforderungsdokument \\
  AP3  & Spezifikation & 1 & \TODO{} & Spezifikationsdokument \\
  AP4  & Feedback Phase 1 einarbeiten & 2 & \TODO{} & überarbeitete Dokumente \\
  AP5  & Architektur und Design & 2 & \TODO{} & Architekturdokument \\
  AP6  & Implementierung & 2 & \TODO{} & Quellcode (GitHub) \\
  AP7  & Feedback Phase 2 einarbeiten & 3 & \TODO{} & überarbeitete Dokumente \\
  AP8  & Testen und Qualitätssicherung & 3 & \TODO{} & Testdokument \\
  AP9  & Benutzeranleitung und Auslieferung & 3 & \TODO{} & ausführbare Anwendung \\
  AP10 & Projektabschluss und Reflexion & 3 & \TODO{} & Abstract, finale Abgabe \\
\end{longtable}
\quelle{Eigene Darstellung.}

\subsection{Meilensteine}
\begin{longtable}{@{}L{1.4cm}L{10.0cm}L{3.8cm}@{}}
  \caption{Meilensteine}\label{tab:meilensteine}\\
  \toprule
  \textbf{ID} & \textbf{Meilenstein} & \textbf{Geplant bis} \\
  \midrule
  \endfirsthead
  \bottomrule
  \endfoot
  M1 & Abgabe Portfolioteil 1 (Konzeptionsphase) & \TODO{Datum} \\
  M2 & Abgabe Portfolioteil 2 (Erarbeitungs- und Reflexionsphase) & \TODO{Datum} \\
  M3 & Funktionsfähige Anwendung, Tests abgeschlossen & \TODO{Datum} \\
  M4 & Abgabe Portfolioteil 3 (Finalisierungsphase) & \TODO{Datum} \\
\end{longtable}
\quelle{Eigene Darstellung.}

\subsection{Zeitplan}
\TODO{Gantt-Diagramm der Arbeitspakete (z. B. mit draw.io) als docs/diagramme/export/zeitplan.pdf}
\begin{figure}[H]
  \caption{Zeitplan der Arbeitspakete}\label{abb:zeitplan}
  \centering
  \grafik{zeitplan}
  \quelle{Eigene Darstellung.}
\end{figure}

\abPhase{2}{% =====================================================================
%  4 Versionsverwaltung  (ab Phase 2)
%  Aufgabenstellung 1.2: „Der Link zum GitHub-Repository ist an geeigneter
%  Stelle in der Projektdokumentation zu ergänzen.“
% =====================================================================
\section{Versionsverwaltung und Repository}\label{sec:versionsverwaltung}
Quellcode, Dokumente, Diagrammquellen und alle Dateien, die zur Ausführung nötig sind, liegen gemeinsam im folgenden Git-Repository auf GitHub:

\url{\GitHubURL}

Das Repository ist öffentlich. Tutor:innen können es daher ohne Einladung und ohne eigenes Konto lesen, klonen oder als ZIP-Datei herunterladen. Die Aufgabenstellung der IU ist urheberrechtlich geschützt und liegt deshalb nicht im Repository. Zusätzlich enthält Anhang~\ref{anh:quellcode} den Quellcode der Anwendung in kopierfähiger Form.

Das Repository ist in vier Hauptordner gegliedert:
\begin{itemize}
  \item \texttt{app/}: die Anwendung mit Quellcode (\texttt{src/balanceos}), Tests (\texttt{tests}), Projektkonfiguration (\texttt{pyproject.toml}) und der Konfiguration für PyInstaller \parencite{pyinstaller}, das daraus die \texttt{.exe} baut. Die ausführbare Datei liegt in \texttt{app/release/BalanceOS.exe}. Die Datei \texttt{app/README.md} beschreibt Installation, Start und Tests.
  \item \texttt{docs/}: die \LaTeX-Quellen aller Dokumente, die TikZ-Quellen der Diagramme mit ihren PDF-Exporten und die gebauten Dokumente in \texttt{docs/pdf}. Der Ordner \texttt{docs/superpowers} enthält die Design- und Planungsnotizen, die bei der Arbeit mit dem \ac{KI}-Werkzeug je Phase entstanden sind.
  \item \texttt{abgabe/}: die eingereichten Abgabe-PDFs und Kurzzusammenfassungen je Phase.
  \item \texttt{tools/}: das Skript \texttt{build\_docs.py}, das Diagramme und Dokumente baut und die Abgabe-PDF einer Phase zusammenstellt.
\end{itemize}

Jede Phase wird auf einem eigenen Zweig bearbeitet, etwa in Phase~3 auf \texttt{phase3}. Am Ende der Phase wird dieser Zweig in \texttt{main} zusammengeführt. Nach der Einreichung markiert ein Tag den eingereichten Stand (\texttt{phase1-abgabe}, \texttt{phase2-abgabe}, \texttt{phase3-abgabe}). So lässt sich jede Abgabe später genau wiederherstellen. Commit-Nachrichten sind auf Deutsch verfasst. In der Implementierung beginnen sie mit dem Inkrement, etwa \enquote{I3: Regelwerk R01--R10 mit Begründung und Evidenzkennzeichnung}. Jedes Inkrement endet mit erfolgreich durchlaufenen Tests und wird auf GitHub hochgeladen. Commits, an denen das \ac{KI}-Werkzeug Claude Code beteiligt war, tragen eine \texttt{Co-Authored-By}-Zeile (siehe Abschnitt~\ref{sec:vorgehen}).
}
\abPhase{3}{% =====================================================================
%  5 Benutzeranleitung  (Phase 3)
%  Aufgabenstellung 1.3: „kurzes Kapitel Benutzeranleitung, welches die
%  Installation und Ausführung der Anwendung beschreibt“.
%  Aufgabenstellung 1: Alle über Standard-Laufzeitumgebungen hinaus
%  benötigten Komponenten explizit beschreiben und mitliefern.
% =====================================================================
\section{Benutzeranleitung}
\TODO{Kurze Hinführung}

\subsection{Systemvoraussetzungen und Installation}
\TODO{Windows 10/11; Python-Version; benötigte Bibliotheken und wie sie mitgeliefert werden; Installation Schritt für Schritt}

\subsection{Start und erste Schritte}
\TODO{Start der Anwendung (z. B. .exe oder python -m …); mitgelieferte Testdaten; typischer Ablauf mit Screenshots}
}
\abPhase{3}{% =====================================================================
%  6 Anforderungserfüllung und technische Schulden  (Phase 3)
%  Aufgabenstellung 1: Nicht vollständig erfüllte funktionale und
%  nichtfunktionale Anforderungen „explizit an geeigneter Stelle als
%  technische Schuld (Technical Debt) dokumentieren“.
%  Aufgabenstellung 1.2: Signifikante Änderungen an Anforderungen und
%  Spezifikation begründen und dokumentieren.
% =====================================================================
\section{Anforderungserfüllung und technische Schulden}
\TODO{Kurze Hinführung}

\subsection{Umsetzungsstand der Anforderungen}
\TODO{Tabelle: Anforderungs-ID, umgesetzt ja/teilweise/nein, Verweis auf Testfall}

\subsection{Technische Schulden}
\TODO{Je Schuld: Beschreibung, Ursache, Auswirkung, geplante Behebung}

\subsection{Änderungen gegenüber der Konzeption}
\TODO{Signifikante Änderungen an Anforderungen und Spezifikation mit Begründung (Feedback der Tutor:innen)}
}

\literaturverzeichnis

\abPhase{2}{%
  \anhang
  \anhangsteil{Quellcode der Anwendung}
  \IfFileExists{build/quellcode.tex}{\input{build/quellcode.tex}}{%
    \TODO{Quellcode-Anhang wird von \texttt{tools/build\_docs.py} erzeugt}}%
}
\end{document}
